% !TeX root = ../main.tex
\chapter{Cohort and Data}
\label{ch:cohort-data}
\thesisplaceholder{Describe the actual study population or biological system, data provenance, and processing.}

\section{Study Setting and Data Provenance}
% 填寫資料來源、版本、收集期間、院所／批次、存取方式與授權條件。
% 非臨床資料改述相應生物系統與取樣設計；不可假造族群或收案程序。

\section{Eligibility, Sampling, and Outcomes}
% 定義納入排除、分析單位、追蹤時間、標籤取得與判定流程。
% 以實際紀錄製作樣本流向及基本特徵；區分人數、檢體數與觀測數。

\section{Quality Control and Preprocessing}
% 記錄重複資料、缺失、異常值、批次、測量範圍、轉換與特徵處理。
% 說明處理步驟在哪個切分上估計；避免以測試資料決定清理或標準化參數。

\section{Ethics, Governance, and Access}
% 填入實際確認的倫理審查／豁免、資料使用、同意與去識別程序。
% 未確認的核准號碼、授權條件或公開資格必須保留待辦，不得推定。
